% Plantilla base: esqueleto con título, subtítulo, leyenda y nota al pie.
% Convención de ids para poder cotejar con fidelidad.md:
%   nodos  \node (n1) ...        aristas  \draw[...] (n1) -- (n2);  % e1
%   totalidades  % g1            ejes/orientaciones  % x1
% Elige el estilo de cada arista por la RELACIÓN que dice el texto (references/lexico.md).
\documentclass[tikz,border=12pt]{standalone}
\usepackage{diagrama}
\begin{document}
\begin{tikzpicture}
  \node[titulo=46] at (4,6.2) {Título: la tesis del texto};
  \node[subtitulo=30] at (4,5.6) {Autor, \textit{Obra}, capítulo o sección};

  \node (n1) at (1.5,4) {Término A};
  \node (n2) at (6.5,4) {Término B\interpretativo};
  \draw[genesis] (n1) -- node[rotulo] {\textit{Nombre de la relación}} (n2);  % e1

  % Leyenda mínima: solo si hay más de tres tipos de trazo y las aristas no llevan su nombre
  \muestra{genesis}{0.2}{0.2}{génesis}
  \muestra{unilateral}{3.4}{0.2}{unilateral}
  \nota{0,-0.5}{\interpretativo{} Lectura del diagramador: lo que el texto no dice con estas palabras.}
\end{tikzpicture}
\end{document}
